\chapter{عنوان پروژه}
یک پروژه در یک فصل نوشته می‌شود و چندین بخش دارد.

\section{کدها}
اگر پروژه شامل کد است، کدها را در پوشه‌ی \lr{codes} قرار داده و سپس به صورت زیر کد خود را در لتک قرار دهید:
\code{python}{codes/example.py}{}

اگر کد در قالب جوپیتر نوتبوک است، ابتدا در جوپیتر نوتبوک خروجی لتک بگیرید سپس فایل لتک را در یک فولدر دلخواه در پوشه‌ی \lr{latex\_chapters} قرار دهید. فایل را کامپایل کنید و pdf نهایی کد را به صورت زیر در متن قرار دهید:
\includepdf[pages=-]{jupyter_codes/code_sample_project/code_project_1.pdf}

\section{تصاویر}
برای وارد کردن تصویر، تصاویر را در پوشه‌ی \lr{figs} قرار دهید. سپس به صورت زیر تصاویر را وارد کنید:
\begin{figure}[h!]
	\centering
	\includegraphics[width=\textwidth]{figs/sample.png}
	\caption{کپشن}
	\label{fig1}
\end{figure}


\section{مراجع}
\lr{bibtex}
مراجع را در \lr{bibs/refs.bib} وارد کنید و سپس به این صورت ارجاع دهید: \cite{kqlstm}.